% GENERATED FILE. Edit canonical lessons, then run npm run generate:books.
% canonical-source-hash: fnv1a64:317a47f61f69eaa4
% canonical-lessons: RU-C352-vina, RU-C352-vinovatyy, RU-C352-vinit, RU-C352-opravdanie, RU-C352-otvetstvennyy, RU-R352-review-345-352

\chapter{Whose fault is it?}
\label{ch:ru-whose-fault-is-it}
\hlchaptermodality{352}

\begin{chapteropening}
\textbf{By the end of this chapter:} \emph{I can say whose fault something is, who is to blame and who is responsible, and talk about an excuse.}

Complete the last lesson of chapter 352: I can say whose fault something is, who is to blame and who is responsible, and talk about an excuse.
\end{chapteropening}

\section[\texorpdfstring{\ru{вина} (viná)}{вина (viná)}]{\ru{вина} (viná) — guilt, fault, blame}

\label{lesson:RU-C352-vina}



\begin{quote}
Before the new one: say the Russian for owing to, then the Russian for therefore.
\end{quote}

\subsection*{You'll want to know: \ru{вина}}
\textbf{\ru{вина}} (\emph{viná}) — \textquotedblleft{}guilt, fault, blame\textquotedblright{}.

\begin{quote}
\ru{Это} \ru{не} \ru{твоя} \ru{вина}.
\end{quote}

\begin{quote}
\emph{Éto ne tvoyá viná.}
\end{quote}

\begin{quote}
It is not your fault.
\end{quote}

Feminine, stressed on the end. \emph{\ru{По} \ru{моей} \ru{вине}} — through my fault. Not to be confused with \emph{\ru{вино}}, wine.

\begin{cousinweb}
\textbf{\ru{вина}} is usually traced to Proto-Slavic \emph{*vina} \textquotedblleft{}guilt, cause\textquotedblright{}; it is the root of \emph{\ru{извиняться}} \textquotedblleft{}to apologise\textquotedblright{}, which asks to be taken out of guilt.
\end{cousinweb}

\begin{grammarlens}[title={What you've built}]
One more word for telling, explaining and describing at B1.
\end{grammarlens}

\subsection*{Your turn}
\begin{itemize}
\raggedright
  \item \emph{Say it:} \emph{viná}
  \item \emph{Say it:} \emph{viná}, once more
  \item \emph{Say it:} a sentence of your own with \emph{viná}
  \item \emph{Recall:} say the Russian for owing to, then the Russian for therefore, then say \emph{viná} again
\end{itemize}

\begin{tcolorbox}[breakable,colback=teal!4,colframe=teal!35!black,title={Before you move on}]
What does \ru{вина} mean? (\textquotedblleft{}Guilt, fault, blame\textquotedblright{}.) Say it once more.
\end{tcolorbox}
\section[\texorpdfstring{\ru{виноватый} (vinovátyy)}{виноватый (vinovátyy)}]{\ru{виноватый} (vinovátyy) — guilty, to blame}

\label{lesson:RU-C352-vinovatyy}



\begin{quote}
Before the new one: say the Russian for therefore, then the Russian for guilt.
\end{quote}

\subsection*{You'll want to know: \ru{виноватый}}
\textbf{\ru{виноватый}} (\emph{vinovátyy}) — \textquotedblleft{}guilty, to blame\textquotedblright{}.

\begin{quote}
\ru{Никто} \ru{не} \ru{виноват}, \ru{так} \ru{получилось}.
\end{quote}

\begin{quote}
\emph{Niktó ne vinovát, tak poluchílos.}
\end{quote}

\begin{quote}
Nobody is to blame; that is how it turned out.
\end{quote}

The short form \emph{\ru{виноват}} is the usual one: \emph{\ru{я} \ru{виноват}} — it is my fault. \emph{\ru{Виноват}!} is also a quick apology.

\begin{cousinweb}
\textbf{\ru{виноватый}} is built on \emph{\ru{вина}} \textquotedblleft{}guilt\textquotedblright{} with the suffix \emph{-\ru{оватый}}; \emph{\ru{вина}} is usually traced to Proto-Slavic \emph{*vina}.
\end{cousinweb}

\begin{grammarlens}[title={What you've built}]
One more word for telling, explaining and describing at B1.
\end{grammarlens}

\subsection*{Your turn}
\begin{itemize}
\raggedright
  \item \emph{Say it:} \emph{vinovátyy}
  \item \emph{Say it:} \emph{vinovátyy}, once more
  \item \emph{Say it:} a sentence of your own with \emph{vinovátyy}
  \item \emph{Recall:} say the Russian for therefore, then the Russian for guilt, then say \emph{vinovátyy} again
\end{itemize}

\begin{tcolorbox}[breakable,colback=teal!4,colframe=teal!35!black,title={Before you move on}]
What does \ru{виноватый} mean? (\textquotedblleft{}Guilty, to blame\textquotedblright{}.) Say it once more.
\end{tcolorbox}
\section[\texorpdfstring{\ru{винить} (vinít)}{винить (vinít)}]{\ru{винить} (vinít) — to blame}

\label{lesson:RU-C352-vinit}



\begin{quote}
Before the new one: say the Russian for guilt, then the Russian for guilty.
\end{quote}

\subsection*{You'll want to know: \ru{винить}}
\textbf{\ru{винить}} (\emph{vinít}) — \textquotedblleft{}to blame\textquotedblright{}.

\begin{quote}
\ru{Она} \ru{винит} \ru{себя} \ru{в} \ru{этой} \ru{аварии}.
\end{quote}

\begin{quote}
\emph{Oná vinít sebyá v étoy avárii.}
\end{quote}

\begin{quote}
She blames herself for this accident.
\end{quote}

Imperfective. \emph{\ru{Винить} \ru{кого} \ru{в} \ru{чём}}: the fault goes after \emph{\ru{в}} + prepositional.

\begin{cousinweb}
\textbf{\ru{винить}} is the verb made from \emph{\ru{вина}} \textquotedblleft{}guilt\textquotedblright{}, which is usually traced to Proto-Slavic \emph{*vina}: to put the guilt on someone.
\end{cousinweb}

\begin{grammarlens}[title={What you've built}]
One more word for telling, explaining and describing at B1.
\end{grammarlens}

\subsection*{Your turn}
\begin{itemize}
\raggedright
  \item \emph{Say it:} \emph{vinít}
  \item \emph{Say it:} \emph{vinít}, once more
  \item \emph{Say it:} a sentence of your own with \emph{vinít}
  \item \emph{Recall:} say the Russian for guilt, then the Russian for guilty, then say \emph{vinít} again
\end{itemize}

\begin{tcolorbox}[breakable,colback=teal!4,colframe=teal!35!black,title={Before you move on}]
What does \ru{винить} mean? (\textquotedblleft{}To blame\textquotedblright{}.) Say it once more.
\end{tcolorbox}
\section[\texorpdfstring{\ru{оправдание} (opravdániye)}{оправдание (opravdániye)}]{\ru{оправдание} (opravdániye) — an excuse, a justification}

\label{lesson:RU-C352-opravdanie}



\begin{quote}
Before the new one: say the Russian for guilty, then the Russian for to blame.
\end{quote}

\subsection*{You'll want to know: \ru{оправдание}}
\textbf{\ru{оправдание}} (\emph{opravdániye}) — \textquotedblleft{}an excuse, a justification\textquotedblright{}.

\begin{quote}
\ru{Для} \ru{такого} \ru{поведения} \ru{нет} \ru{оправдания}.
\end{quote}

\begin{quote}
\emph{Dlya takógo povedéniya net opravdániya.}
\end{quote}

\begin{quote}
There is no excuse for such behaviour.
\end{quote}

Neuter. In court it means an acquittal.

\begin{cousinweb}
\textbf{\ru{оправдание}} is built on \emph{\ru{оправдать}} \textquotedblleft{}to justify\textquotedblright{}, from \emph{\ru{правда}} \textquotedblleft{}truth\textquotedblright{}, inherited from Proto-Slavic \emph{*pravĭda}, built on \emph{*pravŭ} \textquotedblleft{}right, correct\textquotedblright{} (EtymDB). To justify is to show that someone was in the right.
\end{cousinweb}

\begin{grammarlens}[title={What you've built}]
One more word for telling, explaining and describing at B1.
\end{grammarlens}

\subsection*{Your turn}
\begin{itemize}
\raggedright
  \item \emph{Say it:} \emph{opravdániye}
  \item \emph{Say it:} \emph{opravdániye}, once more
  \item \emph{Say it:} a sentence of your own with \emph{opravdániye}
  \item \emph{Recall:} say the Russian for guilty, then the Russian for to blame, then say \emph{opravdániye} again
\end{itemize}

\begin{tcolorbox}[breakable,colback=teal!4,colframe=teal!35!black,title={Before you move on}]
What does \ru{оправдание} mean? (\textquotedblleft{}An excuse, a justification\textquotedblright{}.) Say it once more.
\end{tcolorbox}
\section[\texorpdfstring{\ru{ответственный} (otvétstvennyy)}{ответственный (otvétstvennyy)}]{\ru{ответственный} (otvétstvennyy) — responsible, in charge}

\label{lesson:RU-C352-otvetstvennyy}



\begin{quote}
Before the new one: say the Russian for to blame, then the Russian for an excuse.
\end{quote}

\subsection*{You'll want to know: \ru{ответственный}}
\textbf{\ru{ответственный}} (\emph{otvétstvennyy}) — \textquotedblleft{}responsible, in charge\textquotedblright{}.

\begin{quote}
\ru{Кто} \ru{ответственный} \ru{за} \ru{этот} \ru{проект}?
\end{quote}

\begin{quote}
\emph{Kto otvétstvennyy za étot proyékt?}
\end{quote}

\begin{quote}
Who is responsible for this project?
\end{quote}

\emph{\ru{Ответственный} \ru{за}} + accusative. Short form \emph{\ru{ответствен}}. Of a moment: \emph{\ru{ответственный} \ru{момент}}, a crucial moment.

\begin{cousinweb}
\textbf{\ru{ответственный}} is built on \emph{\ru{ответ}} \textquotedblleft{}an answer\textquotedblright{}, from \emph{\ru{от}-} \textquotedblleft{}back\textquotedblright{} and the old word for counsel and speech, Old Russian \emph{větŭ}. The person responsible is the one who must answer for it.
\end{cousinweb}

\begin{grammarlens}[title={What you've built}]
One more word for telling, explaining and describing at B1.
\end{grammarlens}

\subsection*{Your turn}
\begin{itemize}
\raggedright
  \item \emph{Say it:} \emph{otvétstvennyy}
  \item \emph{Say it:} \emph{otvétstvennyy}, once more
  \item \emph{Say it:} a sentence of your own with \emph{otvétstvennyy}
  \item \emph{Recall:} say the Russian for to blame, then the Russian for an excuse, then say \emph{otvétstvennyy} again
\end{itemize}

\begin{tcolorbox}[breakable,colback=teal!4,colframe=teal!35!black,title={Before you move on}]
What does \ru{ответственный} mean? (\textquotedblleft{}Responsible, in charge\textquotedblright{}.) Say it once more.
\end{tcolorbox}
\section[(dialogue)]{Words from chapters 345–352 — a review}

\label{lesson:RU-R352-review-345-352}



\begin{quote}
Say the Russian for fair and unfair.
\end{quote}

\subsection*{Your turn}
Say these aloud:

\begin{itemize}
\raggedright
  \item fair, unfair, fairness, reasonable, the truth
  \item an advantage, a drawback, benefit, harm, to assess
  \item however, but then, and what is more, so, in other words
  \item an occasion, grounds, a factor, a motive, origin
  \item to contribute to, to lead to, to influence, influence, an impact
  \item a consequence, a result, to affect, an effect, as a result
  \item since, after all, thanks to, owing to, therefore
  \item guilt, guilty, to blame, an excuse, responsible
\end{itemize}

\begin{tcolorbox}[breakable,colback=teal!4,colframe=teal!35!black,title={Before you move on}]
Fair? (\textbf{\ru{справедливый}}.) Unfair? (\textbf{\ru{несправедливый}}.) Fairness? (\textbf{\ru{справедливость}}.) Reasonable? (\textbf{\ru{разумный}}.) The truth? (\textbf{\ru{истина}}.) An advantage? (\textbf{\ru{преимущество}}.) A drawback? (\textbf{\ru{недостаток}}.) Benefit? (\textbf{\ru{польза}}.) Harm? (\textbf{\ru{вред}}.) To assess? (\textbf{\ru{оценивать}}.) However? (\textbf{\ru{однако}}.) But then? (\textbf{\ru{зато}}.) And what is more? (\textbf{\ru{причём}}.) So? (\textbf{\ru{итак}}.) In other words? (\textbf{\ru{иными} \ru{словами}}.) An occasion? (\textbf{\ru{повод}}.) Grounds? (\textbf{\ru{основание}}.) A factor? (\textbf{\ru{фактор}}.) A motive? (\textbf{\ru{мотив}}.) Origin? (\textbf{\ru{происхождение}}.) To contribute to? (\textbf{\ru{способствовать}}.) To lead to? (\textbf{\ru{приводить}}.) To influence? (\textbf{\ru{влиять}}.) Influence? (\textbf{\ru{влияние}}.) An impact? (\textbf{\ru{воздействие}}.) A consequence? (\textbf{\ru{последствие}}.) A result? (\textbf{\ru{следствие}}.) To affect? (\textbf{\ru{отражаться}}.) An effect? (\textbf{\ru{эффект}}.) As a result? (\textbf{\ru{в} \ru{результате}}.) Since? (\textbf{\ru{поскольку}}.) After all? (\textbf{\ru{ведь}}.) Thanks to? (\textbf{\ru{благодаря}}.) Owing to? (\textbf{\ru{вследствие}}.) Therefore? (\textbf{\ru{следовательно}}.) Guilt? (\textbf{\ru{вина}}.) Guilty? (\textbf{\ru{виноватый}}.) To blame? (\textbf{\ru{винить}}.) An excuse? (\textbf{\ru{оправдание}}.) Responsible? (\textbf{\ru{ответственный}}.)
\end{tcolorbox}
